\chapter{Machine timer}
\label{ch:mtimer}

\section{Purpose}

\file{hdl/mtimer.v} is a CLINT-style machine timer. No feature specification
covered it; it was added because the system otherwise has no way to bound a
stalled DMA transfer and no wake source for \sig{WFI}.

Its interrupt is a normal PIC source --- source 7 in the verification environment
--- so it carries no privileged path of its own.

\section{Registers}

\begin{table}[H]
\centering
\caption{Machine timer register map (\file{hdl/mtimer.v})}
\label{tab:tmrmap}
\begin{tabular}{@{}TlllL{0.36\linewidth}@{}}
\toprule
\textnormal{\textbf{Offset}} & \textbf{Register} & \textbf{Access} & \textbf{Reset} & \textbf{Description} \\
\midrule
0x0 & \reg{MTIME\_LO}    & RW & \texttt{0x0000\_0000} & Lower half of the free-running counter. \\
0x4 & \reg{MTIME\_HI}    & RW & \texttt{0x0000\_0000} & Upper half. \\
0x8 & \reg{MTIMECMP\_LO} & RW & \texttt{0xFFFF\_FFFF} & Lower half of the compare value. \\
0xC & \reg{MTIMECMP\_HI} & RW & \texttt{0xFFFF\_FFFF} & Upper half. \\
\bottomrule
\end{tabular}
\end{table}

Any other offset answers SLVERR through the shared slave. Byte strobes are
honoured per lane on all four registers.

\section{Behaviour}

The interrupt is a registered level: \sig{irq\_o $\leftarrow$ (mtime\_q >=
mtimecmp\_q)}. There is no interrupt-status register, and that is deliberate ---
the comparison \emph{is} the status, and the handler clears the line by moving
\reg{mtimecmp} forward. That keeps the timer consistent with the level-sensitive
contract the PIC expects from every peripheral.

\reg{mtimecmp} resets to all ones, so the timer is born disarmed and cannot fire
before software has programmed a deadline.

\subsection{Counter writes}

\sig{mtime} increments every cycle. A software write does not stop it: the
written byte lanes take the software value while the unwritten lanes still take
the increment, so the counter never loses a tick.

\begin{center}
\ttfamily\small
mtime\_q[31:0] <= wr\_time\_lo ? lane\_merge(mtime\_inc[31:0], reg\_wdata, reg\_wstrb)\\
\hspace*{7.4em}: mtime\_inc[31:0];
\end{center}

The upper half is written the same way, so the two halves always advance as a
pair.

\subsection{Arming without a false trigger}

The safe order is \reg{MTIMECMP\_LO} first, then \reg{MTIMECMP\_HI}. While the
upper half is still all ones the comparison cannot pass, so there is no window
between the two 32-bit writes in which a stale low half could fire the interrupt.
Writing the halves in the opposite order does open that window; the software
convention is part of the block's contract.

\begin{figure}[H]
\centering
\begin{tikztimingtable}[timing/slope=0,timing/coldist=2pt,xscale=1.0]
  clk           & 14{C} \\
  MTIMECMP\_LO  & 4D{FFFFFFFF} 10D{target\_lo} \\
  MTIMECMP\_HI  & 8D{FFFFFFFF} 6D{target\_hi} \\
  mtime $\ge$ cmp & 12L 2H \\
  irq           & 14{L} \\
  \extracode
  \begin{pgfonlayer}{background}
    \vertlines[help lines]{4,8,12}
  \end{pgfonlayer}
\end{tikztimingtable}
\caption{Arming sequence. The compare cannot pass while the upper half is all
ones, so writing the low half first is safe. \sig{irq\_o} follows the comparison one
cycle later, since it is registered.}
\label{fig:tmrarm}
\end{figure}

\subsection{Reading the counter from 32-bit software}

\reg{mtime} is architecturally 64 bits and is read through two 32-bit accesses,
so a read can tear: the low half can wrap between the two bus transactions. The
software contract is the standard one --- read \reg{MTIME\_HI}, then
\reg{MTIME\_LO}, then \reg{MTIME\_HI} again, and retry if the high half moved.
No hardware shadow register is provided, because at one count per clock the high
half does not change for roughly $4\times10^{9}$ cycles and the retry costs
nothing on the path that matters.

\section{Use with \sig{WFI}}

Together with \sig{WFI} the timer gives the system an idle mode: program the DMA,
arm the timer as a deadline, execute \sig{WFI}. The CPU stops fetching entirely,
so it generates no interconnect traffic while the DMA works. Either the DMA's
completion interrupt or the timer wakes it. If the timer is what wakes it, the
transfer overran its deadline and the handler can abort the channel.

That sequence is the reason all three delivered blocks exist in one delivery, and
it is exercised end to end in the system testbench.
